\documentclass[12pt,a4paper]{article}

\usepackage{geometry}
\geometry{
    left=2cm, 
    right=2cm,
    top=3cm,  
    bottom=2cm
}

\usepackage[spanish,english]{babel}
\usepackage[utf8]{inputenc}
\usepackage{amsmath}

\usepackage{graphicx}
\usepackage{wrapfig}

\usepackage{csquotes}
\usepackage{hyperref}
\usepackage[style=ieee]{biblatex}
\addbibresource{docs.bib}

\usepackage{setspace}
\setstretch{1.5}

\begin{document}

\section*{SSS - DRD}
\begin{center} 
    \begin{tabular}{|c|l|c|l|}
        \hline
        \textbf{Información} & \multicolumn{3}{|c|}{Software System Specification} \\
        \hline
        \textbf{Proyecto:} & Book Explore & \textbf{Version:} & 10022025 \\    
        \hline
        \textbf{Preparado por:} & Alexis Aguilar & \textbf{Fecha:} & 10/Feb/2025 \\    
        \hline
        \textbf{Recibido por:} & Víctor de la Luz & \textbf{Status:} & En construcción \\    
        \hline
    \end{tabular}
\end{center}

\section{Introduction}
This document aims describing "Book Explorer" system based on exploring new 
titles using a story created by the user.

\section{Applicable and Reference Documents}
This document complies with the definitions contained in the 
ECSS-E-ST-40C standard dated March 6, 2009 \cite{ECSS-E-ST-40C}.

\section{Terms, Definitions and Abbreviated Terms}
\emph{UNAM}: Universidad Nacional Autónoma de México

\section{General Description}
\subsection{Product Perspective}
More than with other systems, its main relationship is with the 
Open Library API, from which most of the information in the books is obtained.
\subsection{General Capabilities}
The system will be able to recommend a series of books to the user using a 
story that the user creates using text displayed on the interface when using 
the system.
\subsection{General Constraints}
None at present.
\subsection{Operational Environment}
Both the system and the user must be connected to the Internet. In general, 
the system makes use of two interfaces: one to interact with the user and 
the other to communicate with Open Library to obtain data from the books.
\subsection{Assumptions and Dependencies}
None at present.

\section{Specific requirements}
\subsection{General}
\subsection{Capabilities Requirements}
\subsection{System interface requirements }

\newpage
\printbibliography

\end{document}